\subsection{Background and relation to earlier work}

Directional zero--one laws are a basic part of the transience theory of
random walks in random environments. Kalikow's work
\cite{Kalikow1981}, together with the regeneration formulation of
Sznitman and Zerner \cite[Lemma~1.1]{SznitmanZerner1999}, established
the sign-union law for uniformly elliptic iid environments,
\[
 P_0(A_\ell\cup A_{-\ell})\in\{0,1\}.
\]
This leaves open which sign is chosen when the union has probability
one. In two dimensions, Zerner and Merkl \cite{ZernerMerkl2001}
proved the full directional zero--one law for iid elliptic environments;
Zerner \cite{Zerner2007} subsequently gave a revised proof. The planar
arguments exploit intersections of paths. The question in higher
dimensions requires a different way to exclude two possible escape signs.

Regeneration is a central tool for this problem and for laws of large
numbers. Sznitman and Zerner
\cite[Theorem~1.4 and Corollary~1.5]{SznitmanZerner1999} constructed
independent increments after suitable record times in an iid environment.
Results on asymptotic directions further distinguish a deterministic
axis from the choice of a sign on that axis: see Simenhaus
\cite{Simenhaus2007} and Drewitz and Ram\'irez
\cite{DrewitzRamirez2010}. A limiting direction or a conditional law of
large numbers does not by itself establish a directional zero--one law.

For dependent environments, Comets and Zeitouni
\cite{CometsZeitouni2004} used environment-independent forced steps to
construct approximate regeneration under mixing assumptions, with
additional hypotheses yielding a law of large numbers. That forcing
device is an antecedent of the marked scripts used here. Rassoul-Agha
\cite{RassoulAgha2005} derived a law of large numbers under Gibbs mixing
assumptions and a directional zero--one hypothesis. Guo
\cite{Guo2014} proved conditional velocity results under a mixing
condition that includes finite-range-dependent environments. These
conclusions concern velocity and do not determine the probabilities of
the two directional escape events. Exact finite-range independence is
also stronger than decay of local correlations: Bramson, Zeitouni, and
Zerner \cite[Theorem~3]{BramsonZeitouniZerner2006} constructed stationary,
uniformly elliptic, polynomially mixing environments in dimension at
least three with positive probability of escape in each of two opposite
directions.

The companion article \cite[Theorem~1.1]{OpenAI2026} proves the iid
directional zero--one law under strict ellipticity: every transition
probability is positive, but no uniform lower bound is required. Here we
extend the uniformly elliptic case to finite-range dependence. The proof
adapts the radius-and-contact strategy developed in
\cite[Sections~2--7]{OpenAI2026}. In that strategy, coexistence first
forces integrability of spatial regeneration radii. Two independent
sequences of regeneration pieces are then placed in an opposed arrangement,
and entropy and endpoint-posterior estimates give incompatible bounds on
their contacts. We retain this architecture and prove all required
estimates for the present environment law. The finite-range extension
depends on keeping track of the rows actually used by a marked path,
replacing intersections by proximity within the dependence range, and
transferring only the stopped path laws justified by separated rows.
No conditional density assumption, independent-label representation,
positive speed, or moment bound on regeneration times is imposed.
